\documentclass[english, parskip=full]{scrartcl}
\usepackage[utf8]{inputenc}  % Kodierung der Datei
\usepackage[english]{babel}  % Sprache auf Deutsch 
\usepackage{color}
\usepackage{graphicx, gensymb}
\usepackage{amsfonts, cancel, amssymb}
\usepackage{amsmath, amsthm, array, esint}
\usepackage{booktabs, multirow}  % schönere Tabellen
\usepackage{caption}  % Anpassung der Captions
\usepackage{scrlayer-scrpage}    % Manipulation des Seitenstils
\pagestyle{scrheadings}  % Stil für die Seite setzen
\automark{section}  % Abschnittsnamen als Seitenbeschriftung 
\setheadsepline{.5pt}    % Kopzeile durch Linie abtrennen
\usepackage[hidelinks]{hyperref}  % Links und weitere PDF-Features
\usepackage{typearea, tikz, pgffor, verbatim}
\usetikzlibrary{calc, arrows.meta,arrows, graphs, quotes, positioning}
\usepackage[cache=false, outputdir=build]{minted}
\usemintedstyle{vs}
\usepackage{placeins} % provides \FloatBarrier

\definecolor{bg}{rgb}{0.9,0.9,0.9}
\newcommand\numberthis{\addtocounter{equation}{1}\tag{\theequation}}
\usepackage{svg}
\usepackage{siunitx}
\usepackage{physics}
\usepackage{tensor}
\renewcommand{\bra}{}
\newcommand{\kom}{}
\newcommand{\brak}{}
\renewcommand{\braket}{}
\newcommand{\intx}{}
\newcommand{\intr}{}
\newcommand{\kla}{}
\newcommand{\klam}{}
\newcommand{\klammer}{}
\newcommand{\oper}{}
\newcommand{\tecxt}{}
\newcommand{\Bra}{}
\newcommand{\Ket}{}
\renewcommand{\i}{\text{i}}
\newcommand{\e}{\text{e}}
\newcommand*{\Fourier}{\textsc{Fourier}\ }
\newcommand*{\F}{\mathcal{F}}
\newcommand*{\sinc}{\text{sinc\,}}
\newcommand{\conv}{\mathop{\scalebox{3.5}{\raisebox{-0.38ex}{\makebox[0.5\width][c]{$\ast$}}}}\limits}

\newcommand*{\titel}{Prime Number Theorem}
\newcommand*{\autor}{Joachim Schwardt}

\hypersetup{pdfauthor={\autor}, pdftitle={\titel}}

%\titlehead{titlehead}
\subject{Analytic Number Theory}
\title{\titel}
\author{\autor}
\date{\today}

\begin{document}

%\begin{titlepage}
\maketitle  % Titel setzen
%\tableofcontents  % Inhaltsverzeichnis setzen
%\end{titlepage}


\section{The prime density function}
Consider the quantity $n!$ for $n\in\mathbb{N}$ and let $p$ be a prime number.
Note that every $p$-th factor of $n!$ is divisible by $p$ once, every $p^2$-th factor is divisible by $p$ twice, and in general every $p^k$-th factor of $n!$ is divisible by $p^k$.
Thus, the number of times $p$ appears in $n!$, which we will call $N_p(n!)$, is given by
\begin{align}
N_p(n!) = \sum_{j\ge 1, jp^1\le n}1 + \sum_{j\ge 1, jp^2\le n}2 + \dots = \sum_{k \ge 1} \sum_{j \ge k, jp^k \le n} k \equiv \sum_{k\ge 1} N_{p,k}(n!). \label{eqn:exact prime factor formula}
\end{align}
Note that $N_{p,k}(n!)$ effectively counts the number of times the factor $p^k$ fits in $n$, which can be written as
\begin{align}
N_{p,k}(n!) &= \left\lfloor \frac{n}{p^k} \right\rfloor.
\end{align}
Thus \autoref{eqn:exact prime factor formula} simplifies to
\begin{align}
N_p(n!) &= \sum_{k \ge 1} \left\lfloor \frac{n}{p^k} \right\rfloor.  \label{eqn:exact prime factor of n factorial}
\end{align}
\autoref{table:exact prime factor example 30fac} illustrates how this formula works for the specific example of $n=30$.
%\begin{table}[!ht]
%\centering
%\begin{tabular}{|*{5}{c|}} \hline
%$p$ & 2 & 3 & 5 & 7 \\ \hline
%$N_{p,1}$ & 5 & 3 & 2 & 1 \\ \hline
%$N_{p,2}$ & 2 & 1 & 0 & 0 \\ \hline
%$N_{p,3}$ & 1 & 0 & 0 & 0 \\ \hline
%$N_{p,k\ge 4}$ & 0 & 0 & 0 & 0 \\ \hline
%$N_p$ & 8 & 4 & 2 & 1 \\ \hline
%\end{tabular}
%\caption{Prime factorization of $10!$ illustrating \autoref{eqn:exact prime factor formula}.}
%\label{table:exact prime factor example}
%\end{table}
%\begin{table}[!ht]
%\centering
%\begin{tabular}{|*{5}{c|}} \hline
%$p$ & 2 & 3 & 5 & 7 \\ \hline
%$N_{p,1}$ & 5 & 3 & 2 & 1 \\ \hline
%$N_{p,2}$ & 2 & 1 & 0 & 0 \\ \hline
%$N_{p,3}$ & 1 & 0 & 0 & 0 \\ \hline
%$N_{p,k\ge 4}$ & 0 & 0 & 0 & 0 \\ \hline
%$N_p$ & 8 & 4 & 2 & 1 \\ \hline
%\end{tabular}
%\begin{tabular}{|*{5}{c|}}\hline
%$p$ & $N_{p,1}$ & $N_{p,2}$ & $N_{p,3}$ & $N_p$ \\ \hline
%2 & 5 & 2 & 1 & 8 \\ \hline
%3 & 3 & 1 & 0 & 4 \\ \hline
%5 & 2 & 0 & 0 & 2 \\ \hline
%7 & 1 & 0 & 0 & 1 \\ \hline
%\end{tabular}
%\caption{Prime factorization of $10!$ illustrating \autoref{eqn:exact prime factor formula}.}
%\label{table:exact prime factor example 10fac}
%\end{table}
\begin{table}[!ht]
\centering
\begin{tabular}{|*{6}{c|}}\hline
$p$ & $N_{p,1}$ & $N_{p,2}$ & $N_{p,3}$ & $N_{p,4}$ & $N_p$ \\ \hline
2 & 15 & 7 & 3 & 1 & 26 \\ \hline
3 & 10 & 3 & 1 & 0 & 14 \\ \hline
5 & 6 & 1 & 0 & 0 & 7 \\ \hline
7 & 4 & 0 & 0 & 0 & 4 \\ \hline
11 & 2 & 0 & 0 & 0 & 2 \\ \hline
13 & 2 & 0 & 0 & 0 & 2 \\ \hline
17 & 1 & 0 & 0 & 0 & 1 \\ \hline
19 & 1 & 0 & 0 & 0 & 1 \\ \hline
23 & 1 & 0 & 0 & 0 & 1 \\ \hline
29 & 1 & 0 & 0 & 0 & 1 \\ \hline
\end{tabular}
\caption{Prime factorization of $30!$ illustrating \autoref{eqn:exact prime factor of n factorial}.}
\label{table:exact prime factor example 30fac}
\end{table}
The floor function is connected to the fractional part via $\lfloor x \rfloor = x - \{ x \}$, which we may use to rewrite the counting function as
\begin{align}
N_p(n!) &= \sum_{k \ge 1} \frac{n}{p^k} - \sum_{k\ge 1} \left\{ \frac{n}{p^k} \right\} = \frac{n}{p - 1} - \sum_{k\ge 1} \left\{ \frac{n}{p^k} \right\}.
\end{align}
%Using this equation, we can rewrite the factorial as a product over the primes
%\begin{align}
%%n! &= \prod_p p^{N_p(n!)} = \prod_p \prod_{k\ge 1} p^{\left\lfloor \frac{n}{p^k} \right\rfloor}.
%n! &= \prod_p p^{N_p(n!)} = \prod_p p^{\frac{n}{p-1}} \prod_{k\ge 1} p^{-\left\{ \frac{n}{p^k} \right\}}.
%\end{align}
%We take the logarithm to convert products into sums, leading to
%\begin{align}
%\log n! &= \sum_p \klam\left[ \frac{n}{p-1}\log p - \sum_{k\ge 1} \klammer\left\{ \frac{n}{p^k} \right\} \log p \right] = \\
%&= \sum_p \log p \klam\left[ \frac{n}{p-1} - \sum_{k\ge 1} \klammer\left\{ \frac{n}{p^k} \right\} \right].
%\end{align}
Using this equation, we can rewrite the factorial as a product over the primes
\begin{align}
n! &= \prod_p p^{N_p(n!)}.
\end{align}
We take the logarithm to convert products into sums, leading to
\begin{align}
\log n! &= \sum_p N_p(n!) \log p = \sum_{p\le n} \log p \klam\left[ \frac{n}{p-1} - \sum_{k\ge 1} \klammer\left\{ \frac{n}{p^k} \right\} \right].
\end{align}
Now let us turn to a different approximation for $n!$ using the integral
\begin{align}
n! &= \intx\int_{0}^{\infty}x^n \e^{-x}\, \mathrm{d}x = \intx\int_{0}^{\infty} \e^{-x + n\log x}\, \mathrm{d}x.
\end{align}
The function $-x + n\log x$ has a minimum at $x = n$.
Expanding around this point gives
\begin{align}
-x + n\log x &\sim -n + n\log n - \frac{1}{2n} \kla\left( x - n \right)^2
\end{align}
This leads to the approximation
\begin{align}
n! &\sim \intx\int_{0}^{\infty} \e^{-n + n\log n} \e^{-\frac{1}{2n} (x-n)^2} \, \mathrm{d}x \kom\overset{\substack{{z\, := \, \frac{x - n}{\sqrt{n}}} \\ |}}{=} \e^{-n + n\log n} \intx\int_{-\sqrt{n}}^{\infty} \e^{-\frac{1}{2}z^2} \, \sqrt{n}\mathrm{d}z.
\end{align}
We can compute the upper part of the integral analytically, but need to make further approximations for the lower part.
%In a rather drastic attempt, we will simply use the midpoint rule with a single support point at $x=-\frac{1}{2n}$, giving
%\begin{align}
%\intx\int_{-\frac{1}{n}}^{0} \e^{-\frac{n}{2}x^2} \, \mathrm{d}x &\sim \e^{-\frac{1}{8n^3}}
%\end{align}
Since we are interested in the behaviour for large $n$, we consider the small deviation from a gaussian integral
\begin{align}
\epsilon_n &:= \intx\int_{-\infty}^{-\sqrt{n}} \e^{-\frac{1}{2}x^2} \, \mathrm{d}x \kom\overset{\substack{{z\, :=\, \e^{-\frac{1}{2}x^2}} \\ |}}{=} \intx\int_{0}^{\e^{-\frac{n}{2}}} z \, \frac{\mathrm{d}z}{z \sqrt{-2\log z}}.
\end{align}
Expanding around the upper bound yields
\begin{align}
\epsilon_n &\sim \intx\int_{0}^{\e^{-\frac{n}{2}}} \frac{1}{\sqrt{n}} - \frac{1}{2\sqrt{n}^3} \frac{-2}{\e^{-\frac{n}{2}}} \kla\left( z - \e^{-\frac{n}{2}} \right) \, \mathrm{d}z = \frac{\e^{-\frac{n}{2}}}{\sqrt{n}} - \frac{\e^{-\frac{n}{2}}}{2\sqrt{n^3}} .
\end{align}
Collecting all the terms, we arrive at
\begin{align}
n! &\sim \e^{-n + n\log n} \klam\left[ \sqrt{2\pi n} - \epsilon_n \right] \sim \e^{-n + n\log n} \klam\left[ \sqrt{2\pi n} - \frac{\e^{-\frac{n}{2}}}{\sqrt{n}} \right].
\end{align}
Taking the logarithm of this expression gives
\begin{align}
\log n! &\sim -n + n\log n + \log \kla\left( \sqrt{2\pi n} - \frac{1}{\sqrt{n}} \e^{-\frac{n}{2}} \right) \sim \\
&\sim n\log n - n + \frac{\log n}{2}.
\end{align}
Comparing to our previous result, we find
\begin{align}
\sum_{p\le x} \log p \klam\left[ \frac{x}{p-1} - \sum_{k\ge 1} \klammer\left\{ \frac{x}{p^k} \right\} \right] &\sim x\log x - x + \frac{\log x}{2},\\
\Rightarrow\ \ \ \sum_{p\le x} \log p \klam\left[ \frac{1}{p-1} - \sum_{k\ge 1} \frac{1}{x} \klammer\left\{ \frac{x}{p^k} \right\} \right] &\sim \log x - 1 + \frac{\log x}{2x}.
\end{align}
Consider some expression of the form
\begin{align}
F(x) &:= \sum_{p\le x} f(p).
\end{align}
Differentiating this with respect to $x$ is not possible, since this is a step-function.
However, for small $\Delta x<1$, the difference
\begin{align}
F(x+\Delta x) - F(x)  &= \sum_{p-x\le \Delta x} f(p) = \begin{cases} f(x) & x\ \tecxt\text{prime},\\ 0 \end{cases}
\end{align}
gives us a notion of ``local change''.
Note that this function encapsulates the prime density $\delta(x)$, in the sense that this difference amounts to $f(x)$ for the fraction of numbers near $x$ that are prime.
Instead of the step function, we may therefore approximate the ``derivative'' as $\delta(x)f(x)$.
In this way, we arrive at
\begin{align}
\delta(x) \frac{\log x}{x-1} &\sim \frac{1}{x}
\end{align}
where we now neglected the correction from the fractional part and the correction to the analytic approximation of $\log x!$.
Thus the density of primes is given by
\begin{align}
\delta (x) &\sim \frac{1 - \frac{1}{x}}{\log x} .
\end{align}
The prime counting function can therefore be approximated by the integral
\begin{align}
\pi(x) &\sim \intx\int_{2}^{x+1} \delta(x) \, \mathrm{d}x \sim \intx\int_{2}^{x+1} \frac{1-\frac{1}{t}}{\log t} \, \mathrm{d}t.
\end{align}
Numerical analysis shows that the correction $-\frac{1}{t}$ can safely be neglected, as it does not contribute to the relative error for large $x$.
\\
Let us try to improve this estimate by inspecting the fractional error term
\begin{align}
\epsilon(x) &:= \sum_{p\le x} \log p \sum_{k\ge 1} \klammer\left\{ \frac{x}{p^k} \right\}.
\end{align}
%For $x \gg p^k$, the expression $\klammer\left\{ \frac{x}{p^k} \right\}$ can be very roughly interpreted as random noise.
%Let $\epsilon \ll 1$.
%Then $\frac{x}{(\epsilon x)^2} < 1$ for $x > \frac{1}{\epsilon^2}$ which allows us to write
%\begin{align}
%\sum_{p\le x} \sum_{k\ge 1} \klammer\left\{ \frac{x}{p^k} \right\} &= \sum_{\epsilon x< p \le x} \sum_{k\ge 1} \klammer\left\{ \frac{x}{p^k} \right\} + \sum_{p\le \epsilon x} \sum_{k\ge 1} \klammer\left\{ \frac{x}{p^k} \right\} = \\
%&= \sum_{\epsilon x< p \le x} \klammer\left\{ \frac{x}{p} \right\} + \sum_{\epsilon x< p \le x} \sum_{k\ge 2} \frac{x}{p^k} + \sum_{p\le \epsilon x} \klammer\left\{ \frac{x}{p} \right\} + \sum_{p\le \epsilon x} \sum_{k\ge 2} \klammer\left\{ \frac{x}{p^k} \right\}.
%\end{align}
%Note that the first and third part may be combined and roughly approximated as $\pi(x) \bra\langle \klammer\left\{ \frac{x}{p} \right\} \rangle$, where the expectation value of the fractional part is about $\frac{1}{2}$.
%This leads to
%\begin{align}
%\dots &\approx \frac{\pi(x)}{2} + x\sum_{\epsilon x<p\le x} \frac{1}{p^2} \frac{1}{1 - \frac{1}{p}} + \sum_{p\le \epsilon x} \sum_{k\ge 2} \klammer\left\{ \frac{x}{p^k} \right\}.
%\end{align}
%For very small $\epsilon$, we may simplify this expression to
%\begin{align}
%\sum_{p\le x} \sum_{k\ge 1} \klammer\left\{ \frac{x}{p^k} \right\} &\approx \frac{\pi(x)}{2} + x\sum_{p\le x} \frac{1}{p(p-1)}.
%\end{align}
For $x \gg p^k$, the expression $\klammer\left\{ \frac{x}{p^k} \right\}$ can be very roughly interpreted as random noise.
%Let $\epsilon \ll 1$.
%Then $\frac{x}{(\epsilon x)^2} < 1$ for $x > \frac{1}{\epsilon^2}$.
%For $p > \epsilon x$ we may then write
%\begin{align}
%\sum_{k\ge 1} \klammer\left\{ \frac{x}{p^k} \right\} &= \klammer\left\{ \frac{x}{p} \right\} + \sum_{k\ge 2} \frac{x}{p^k}.
%\end{align}
%Note that on average the first part is just going to be about $\frac{1}{2}$.
%This leads to
%\begin{align}
%\sum_{k\ge 1} \klammer\left\{ \frac{x}{p^k} \right\} &\approx \frac{1}{2} + \frac{x}{p(p-1)}.
%\end{align}
%Inserting into our previous identity, we now get
%\begin{align}
%\sum_{p\le x} \log p \klam\left[ \frac{1}{p-1} - \frac{1}{2x} - \frac{1}{p(p-1)} \right] &\approx \log x - 1 + \frac{\log x}{2x}.
%\end{align}
%Note that within the approximation $\delta(x) \approx \frac{1}{\log x}$, we may use an integral to evaluate $\sum_{p\le x} \log p \approx x$.
%Thus the left hand side may be rewritten as
%\begin{align}
%-\frac{1}{2} + \sum_{p\le x} \frac{\log p}{p} &\approx \log x - 1.
%\end{align}
%This leads to $\delta(x) \approx \frac{1}{\log x}$, which we already found before, but now the correction term disappeared.
Note that
\begin{align}
\sum_{k\ge 1} \klammer\left\{ \frac{x}{p^k} \right\} &= \sum_{k = 1}^{\kappa-1} \klammer\left\{ \frac{x}{p^k} \right\} + \sum_{k\ge \kappa} \frac{x}{p^k},
\end{align}
where $\kappa$ is chosen to satisfy the condition $\frac{x}{p^\kappa} < 1$.
This can be rewritten as $\kappa > \frac{\log x}{\log p}$.
For the majority of primes, this cut-off will be as small as $\kappa=2$.
More precisely, the condition is $p > \sqrt{x}$.
Thus we find
\begin{align}
\epsilon(x) &= \sum_{p\le \sqrt{x}} \log p \sum_{k\ge 1} \klammer\left\{ \frac{x}{p^k} \right\} + \sum_{\sqrt{x} < p \le x} \log p \klam\left[ \klammer\left\{ \frac{x}{p} \right\} + \sum_{k\ge 2} \frac{x}{p^k} \right] = \\
&= \sum_{p\le \sqrt{x}} \log p \klam\left[ \sum_{k=1}^{\kappa-1} \klammer\left\{ \frac{x}{p^k} \right\} + \frac{x}{p^{\kappa-1}(p-1)} \right] + \sum_{\sqrt{x} < p \le x} \log p\klam\left[ \klammer\left\{ \frac{x}{p} \right\} + \frac{x}{p(p-1)} \right].
\end{align}
So far we have not approximated anything.
The sum containing $\klammer\left\{ \frac{x}{p} \right\}$ approaches some constant numerical value slightly below $\frac{1}{2}$\footnote{For $x=10^8$ we find $\sum_{\sqrt{x} < p \le x} \log p\klammer\left\{ \frac{x}{p} \right\} = 0.4226...$ and $\sum_{p \le x} \log p \klammer\left\{ \frac{x}{p} \right\} = 0.4227...$. 
This is consistent with $\gamma - 1$, where $\gamma = 0.577...$ is the Euler-Mascheroni constant.%http://www.ebyte.it/library/educards/constants/MathConstants.html
}.
Since it is not too far off, and the term will have order $\frac{1}{x}$, we will continue with the rough approximation of $\frac{1}{2}$ regardless.
Similarly, the first fractional term is many orders of magnitude smaller than the one we just discussed, and will therefore simply be ignored.
It turns out that the geometrical part of this sum is even smaller still, so we are left with
\begin{align}
\epsilon(x) &\approx \frac{1}{2} + x \sum_{\sqrt{x} < p \le x} \frac{\log p}{p(p-1)}.
\end{align}
NEW ATTEMPT:
\\
We know that
\begin{align}
\log x! &= \sum_{p\le x} \log p\klam\left[ \sum_{k\ge 1} \frac{x}{p^k} - \sum_{k\ge 1} \left\{ \frac{x}{p^k} \right\} \right] = \\
&= \sum_{\sqrt{x} < p\le x} \log p\klam\left[ \frac{x}{p} - \left\{ \frac{x}{p} \right\} \right] + \sum_{p\le \sqrt{x}} \log p\klam\left[ \frac{x}{p-1} - \sum_{k\ge 1} \left\{ \frac{x}{p^k} \right\} \right]. %= \\
%&= \sum_{\sqrt{x} < p \le x} \log p \left\lfloor \frac{x}{p} \right\rfloor
\end{align}
The first fractional expression appears to converge to $(\gamma-1) x$, although we do not see how to provide a proof of this claim.
Similarly, the second fractional expression appears to converge to $-\sqrt{x}$ with very little error.
Since these are just correction terms, we will continue with their approximate forms.
Then we find
\begin{align}
\log x! &\approx (\gamma - 1)x - \sqrt{x} + x\sum_{\sqrt{x} < p\le x} \frac{\log p}{p} + x\sum_{p\le \sqrt{x}} \frac{\log p}{p-1}.
\end{align}
Note that the square root term is negligible compared to the rest and that one can not approximate the $-1$ in the denominator of the last term, as that introduces a relatively large error.
However, we may introduce this factor in the other summation, since it only contributes of the order of $\sqrt{x}$ in that case.
This leaves us with
\begin{align}
\log x! &\approx (\gamma-1) x + x \sum_{p\le x} \frac{\log p}{p-1}.
\end{align}
We may implicitly define the prime density function $\delta(x)$ via
\begin{align}
\sum_{p\le x} f(p) &\overset{!}{\sim} \intx\int_{2}^{x+1} \delta(t) f(t)\, \mathrm{d}t.
\end{align}
Doing so allows us to rewrite the previous equation as
\begin{align}
x\log x - x &\sim (\gamma-1) x + x \intx\int_{2}^{x+1}\delta(t) \frac{\log t}{t-1}\, \mathrm{d}t,\\
\Rightarrow\ \ \ \log x &\sim \gamma + \intx\int_{2}^{x+1} \delta(t) \frac{\log t}{t-1}\, \mathrm{d}t.
\end{align}
There appears to be a mystery constant $\alpha \approx 0.833...$ such that
\begin{align}
\delta(x) &\approx \frac{1 - \frac{\alpha}{x}}{\log x}.
\end{align}



%Define
%\begin{align}
%E(x) &= \sum_{p\le x} e(p,x)
%\end{align}
%and note that
%\begin{align}
%E(x+\Delta x) - E(x) &= \sum_{p\le x+\Delta x} e(p,x+\Delta x) - \sum_{p\le x} e(p,x) \sim \\
%&\sim \sum_{p-x\le \Delta x} e(p,x) + \Delta x\sum_{p\le x+\Delta x} \partial_xe(p,x)\sim \\
%&\sim \Delta x\delta(x) e(x,x) + \Delta x\sum_{p\le x}\partial_x e(p, x).
%\end{align}




%\begin{thebibliography}{9}
%\bibitem{example} description, \url{https://www.exmapleurl}, day.\,month~2021.
%\end{thebibliography}

\end{document}